% =============================================================
%  Asymmetric Numeral Systems (ANS) — Improved Exposition
% =============================================================

\section{Asymmetric Numeral Systems (ANS)}

% -----------------------------------------------------------
\subsection{Why ANS Was Invented}
% -----------------------------------------------------------

To compress data, we want to assign short bit-strings to
\emph{common} symbols and longer bit-strings to \emph{rare}
symbols.  Shannon's entropy theorem tells us the ideal length
for a symbol with probability $p$ is $-\log_2 p$ bits.  Two
classical methods try to achieve this:

\begin{itemize}[leftmargin=2em]
    \item \textbf{Huffman coding} assigns a fixed binary
          codeword to each symbol.  It is very fast, but
          codeword lengths must be whole numbers, so it wastes
          bits whenever the ideal length $-\log_2 p$ is not an
          integer.  For example, if $p = 0.3$ the ideal length
          is $\approx 1.74$ bits, but Huffman is forced to use
          either 1 or 2 bits.

    \item \textbf{Arithmetic coding} avoids this waste by
          encoding the entire message as a single fractional
          number, allowing it to approach the entropy limit
          arbitrarily closely.  The drawback is that it
          requires expensive division operations and careful
          multi-precision arithmetic.
\end{itemize}

\medskip
In 2009, Jarek Duda introduced \textbf{Asymmetric Numeral
Systems (ANS)}: an entropy coder that achieves the compression
quality of arithmetic coding while running nearly as fast as
Huffman coding.  ANS is now used in
\textbf{Zstandard (zstd)}, \textbf{LZFSE} (Apple), and other
widely deployed compressors.

\medskip
\textbf{Road-map.}  This section builds up ANS in three
stages.
\begin{enumerate}
  \item We explain the core idea: a single integer $x$ that
        carries all compressed information.
  \item We develop \textbf{rANS} (Range ANS), whose
        encode/decode steps are explicit arithmetic formulas.
        This is the easiest variant to understand.
  \item We show how \textbf{tANS} (Tabled ANS) replaces
        that arithmetic with a pre-built lookup table,
        eliminating every division at runtime.
\end{enumerate}

% -----------------------------------------------------------
\subsection{The Central Idea: One Integer Carries Everything}
% -----------------------------------------------------------

Most compressors build up a bitstream symbol by symbol.  ANS
takes a different approach: it maintains a single integer
$x$ — called the \textbf{state} — and encodes the entire
message by repeatedly updating $x$.  At the end, $x$ itself,
written in binary, \emph{is} the compressed output.

\begin{definitionbox}
\textbf{State $x$.}  The encoder starts from an agreed
initial value (e.g.\ $x = L$, defined below).  Each time a
symbol is encoded, $x$ is replaced by a new, larger integer
that ``contains'' both the old $x$ and the new symbol.  When
all symbols have been encoded, the final $x$ is written out
as the compressed bitstream.  The decoder reads that integer,
and by reversing the updates recovers the original symbols
one at a time.
\end{definitionbox}

\subsubsection*{Why does this work as compression?}

Each update multiplies $x$ by roughly $1/p_s$, where $p_s$
is the probability of the symbol just encoded.  That adds
$\log_2(1/p_s) = -\log_2 p_s$ bits to the binary
representation of $x$ --- exactly the number of bits that
information theory says the symbol deserves.  Common symbols
(large $p_s$) grow $x$ only a little; rare symbols (small
$p_s$) grow $x$ a lot.  After encoding $n$ symbols the size
of $x$ in bits closely tracks the Shannon entropy of the
sequence.

% -----------------------------------------------------------
\subsection{Setting Up: Symbols, Frequencies, and the Table Size}
% -----------------------------------------------------------

Before encoding anything, we need to fix a few parameters.
The table below defines every variable used throughout this
section; refer back to it whenever a symbol is unfamiliar.

\begin{center}
\renewcommand{\arraystretch}{1.3}
\begin{tabular}{@{}lp{10cm}@{}}
\toprule
\textbf{Variable} & \textbf{Meaning} \\
\midrule
$\mathcal{A} = \{s_1,\ldots,s_k\}$
  & The \textbf{alphabet}: the set of symbols to be compressed. \\
$p_i$
  & The \textbf{true probability} of symbol $s_i$; satisfies
    $\sum_i p_i = 1$. \\
$R$
  & A positive integer (e.g.\ 8, 12, or 16) that controls
    precision. \\
$L = 2^R$
  & The \textbf{table size}: the total number of slots
    available to the probability model. \\
$f_i$
  & The \textbf{frequency} of symbol $s_i$: a positive integer
    approximating $p_i \cdot L$.  Must satisfy $\sum_i f_i = L$. \\
$\tilde{p}_i = f_i / L$
  & The \textbf{quantised probability} of $s_i$: the
    approximation to $p_i$ that ANS actually uses. \\
$c_i = \sum_{j < i} f_j$
  & The \textbf{cumulative frequency} of $s_i$: the sum of
    all frequencies of symbols that come \emph{before} $s_i$
    in the alphabet ordering.  Satisfies $c_1 = 0$. \\
$x$
  & The current encoder/decoder \textbf{state}: a single
    non-negative integer. \\
$[L,\,2L)$
  & The \textbf{normalised range}: the interval $\{L, L{+}1,
    \ldots, 2L{-}1\}$ in which $x$ is always kept. \\
\bottomrule
\end{tabular}
\end{center}

\subsubsection*{Why integers instead of exact probabilities?}

Exact probabilities are usually irrational numbers that cannot
be stored or computed efficiently.  By rounding each $p_i$ to
the nearest multiple of $1/L$, we get frequencies $f_i$ that
are small whole numbers, and all arithmetic reduces to
integer operations.  Choosing $L = 2^R$ is especially
convenient because division and remainder by $L$ become
single bit-shift and bit-mask instructions on any modern CPU.

\subsubsection*{What is $c_i$ (cumulative frequency) for?}

Think of the $L$ slots as a row of pigeonholes numbered $0$
to $L-1$.  Symbol $s_i$ is assigned a contiguous block of
$f_i$ pigeonholes, starting at pigeonhole $c_i$ and ending
just before $c_i + f_i$.  Because the blocks are contiguous
and non-overlapping, the decoder can identify any symbol
simply by checking which block a number falls into.


\begin{examplebox}
\textbf{Example: assigning slot ranges (with explicit pigeonholes).}
Suppose $\mathcal{A} = \{A, B, C\}$, $L = 8$, and the
frequencies are $f_A = 4$, $f_B = 2$, $f_C = 2$
(so $4+2+2=8=L$, as required).

\medskip
\begin{center}
\begin{tabular}{@{}lllll@{}}
\toprule
Symbol & Frequency $f$ & Cumulative $c$ & Prob.\ $\tilde{p}=f/L$
  & Pigeonhole numbers \\
\midrule
$A$ & 4 & $c_A=0$         & 0.50 & $0, 1, 2, 3$ \\
$B$ & 2 & $c_B=0+4=4$     & 0.25 & $4, 5$ \\
$C$ & 2 & $c_C=0+4+2=6$   & 0.25 & $6, 7$ \\
\bottomrule
\end{tabular}
\end{center}

\medskip
\noindent
The pigeonholes tile $[0,8)$ perfectly, with each symbol
occupying a contiguous block:

\medskip
\begin{center}
\ttfamily
\begin{tabular}{ccccccccc}
\toprule
Pigeonhole \#: & 0 & 1 & 2 & 3 & 4 & 5 & 6 & 7 \\
\midrule
Symbol owner:  & A & A & A & A & B & B & C & C \\
\bottomrule
\end{tabular}
\end{center}

\medskip
\noindent
The decoder computes $r = x \bmod 8$.  
If $r = 5$, it falls in pigeonhole $5$ $\rightarrow$ symbol $B$.  
If $r = 2$, it falls in pigeonhole $2$ $\rightarrow$ symbol $A$.  
If $r = 7$, it falls in pigeonhole $7$ $\rightarrow$ symbol $C$.  
This is exactly how the decoder identifies symbols from the state $x$.
\end{examplebox}


% -----------------------------------------------------------
\subsection{Two Variants: rANS and tANS}
% -----------------------------------------------------------

\begin{center}
\begin{tabular}{@{}llll@{}}
\toprule
\textbf{Variant} & \textbf{Full name} & \textbf{Encode/decode via}
  & \textbf{Used in} \\
\midrule
\textbf{rANS} & Range ANS   & Arithmetic formulas & Zstandard, CRAM \\
\textbf{tANS} & Tabled  ANS & Lookup tables       & Zstandard, LZFSE \\
\bottomrule
\end{tabular}
\end{center}

Both variants share the same mathematical idea.  We study
rANS first because its formulas are easy to inspect and
verify.  Section~\ref{sec:tans} then shows how tANS
pre-computes those formulas into two arrays so that no
arithmetic is needed at runtime.

% -----------------------------------------------------------
\subsection{rANS: Range Asymmetric Numeral Systems}
% -----------------------------------------------------------

\subsubsection{Keeping the State Bounded}

If we simply multiplied $x$ by $1/p_s$ after every symbol,
$x$ would grow without bound, requiring ever more memory and
time to manipulate.  We prevent this by keeping $x$ inside a
fixed \textbf{normalised range} $[L, 2L)$ at all times.

Whenever $x$ grows too large, we peel off its lowest-order
bit and write it to the bitstream (this is called
\textbf{emitting} a bit).  Whenever $x$ shrinks too small,
we read the next bit from the bitstream and glue it onto $x$
from below (this is called \textbf{renormalising upward}).
These bits are the compressed output; the final value of $x$
is appended at the end so the decoder knows where to start.


\subsubsection{The rANS Encoding Step}
\label{sec:rans_encode}

\paragraph{What we want to achieve.}
Given the current state $x$ and a symbol $s$, we want to
produce a new state $x'$ such that:
\begin{enumerate}
  \item $x'$ encodes the information ``symbol $s$ was seen''.
  \item The decoder can recover both $s$ \emph{and} the
        original $x$ from $x'$ alone.
  \item $x'$ stays in the normalised range $[L, 2L)$.
\end{enumerate}

\paragraph{The encoding formula.}
The formula that achieves all three goals is:
\begin{equation}
  \label{eq:rans_enc}
  x' \;=\; \Bigl\lfloor \frac{x}{f_s} \Bigr\rfloor \cdot L
          \;+\; c_s \;+\; (x \bmod f_s)
\end{equation}
where $f_s$ is the frequency of $s$ and $c_s$ is its
cumulative frequency.

Let's break down why this formula works.

\paragraph{Why goal 1 holds (information content).}
The term $\lfloor x/f_s \rfloor \cdot L$ multiplies $x$ by
approximately $L/f_s$.  Since $L/f_s = 1/\tilde{p}_s$, this
increases the size of $x$ by a factor of $1/\tilde{p}_s$.
In bit terms, we add $\log_2(1/\tilde{p}_s) = -\log_2 \tilde{p}_s$
bits — exactly the ideal information content of symbol $s$.

Common symbols (large $f_s$) cause a small increase; rare
symbols (small $f_s$) cause a large increase.

\paragraph{Why goal 2 holds (decodability).}
The decoder will identify $s$ by looking at $x' \bmod L$.
Let's compute it:
\[
x' \bmod L \;=\; c_s + (x \bmod f_s)
\]

Two key facts about the right-hand side:
\begin{itemize}
  \item $0 \le x \bmod f_s < f_s$
  \item Therefore $c_s \le x' \bmod L < c_s + f_s$
\end{itemize}

But $[c_s,\; c_s+f_s)$ is precisely the slot range that
belongs to symbol $s$!  So:

\begin{center}
\boxed{x' \bmod L \text{ always falls inside } s\text{'s pigeonhole block.}}
\end{center}

The decoder simply computes $r = x' \bmod L$, finds which
symbol owns that pigeonhole, and recovers $s$ instantly.

\paragraph{Why goal 3 holds (staying in $[L, 2L)$).}
The formula above might produce an $x'$ that is too large.
To prevent this, we \textbf{pre-normalise} $x$ before applying
the formula.

\textbf{The trick:} Before encoding, we repeatedly shrink $x$
by shifting out its lowest bits (and writing those bits to the
output) until:
\[
f_s \;\le\; x \;<\; 2f_s
\]

Why $[f_s, 2f_s)$?  Because when $x$ is in this range, we have
$\lfloor x/f_s \rfloor = 1$.  Then the encoding formula
simplifies dramatically:
\[
x' = 1 \cdot L + c_s + (x \bmod f_s) = L + c_s + (x \bmod f_s)
\]

Now check the range of $x'$:
\begin{itemize}
  \item Minimum: when $x \bmod f_s = 0$ → $x' = L + c_s \ge L$
  \item Maximum: when $x \bmod f_s = f_s-1$ → $x' = L + c_s + (f_s-1) < L + c_s + f_s$
\end{itemize}

Since $c_s + f_s \le L$ (the pigeonhole blocks fit exactly in
$[0, L)$), we get $x' < L + L = 2L$.  Thus $x' \in [L, 2L)$.

\paragraph{Why pre-normalise before, not after?}
The low bits of $x'$ (specifically $x' \bmod L$) contain the
symbol tag that the decoder needs.  If we applied the formula
first and then tried to shrink $x'$ by shifting, we would
destroy that tag.  By shrinking $x$ \emph{before} the formula,
the tag is freshly created in the correct position and never
needs to be modified.

\begin{importantbox}
\textbf{Analogy: a water tank.}
Think of $x$ as a water tank.  Encoding a symbol adds an
amount of water proportional to $-\log_2 p_s$ (rare symbols
add more).  When the tank gets too full ($x \ge 2f_s$), we
pour out the excess (emit bits) \emph{before} adding more
water.  This keeps the tank level between $f_s$ and $2f_s$
before each addition, ensuring the final level stays between
$L$ and $2L$.
\end{importantbox}

\begin{definitionbox}
\textbf{rANS Encode Algorithm.}  To encode symbol $s$
(with frequency $f_s$ and cumulative frequency $c_s$)
into the current state $x$:

\medskip
\textbf{Step 1 — Pre-normalise (shrink $x$ until it fits):}
\begin{algorithmic}[1]
\While{$x \ge 2f_s$}
    \State $\text{emit}(x \mathrel{\&} 1)$
           \Comment{write the lowest bit of $x$ to the output}
    \State $x \leftarrow x \gg 1$
           \Comment{right-shift (divide by 2)}
\EndWhile
\Comment{Now $f_s \le x < 2f_s$}
\end{algorithmic}

\medskip
\textbf{Step 2 — Encode (apply the simplified formula):}
\begin{algorithmic}[1]
\State $x \leftarrow L + c_s + (x \bmod f_s)$
       \Comment{new state, guaranteed in $[L, 2L)$}
\end{algorithmic}
\end{definitionbox}

\begin{importantbox}
\textbf{Encoding order and reversal.}
Notice that bits are emitted \emph{before} the formula is
applied.  This means the bits corresponding to the \emph{last}
symbol encoded appear first in the output stream.  Consequently,
the decoder will recover symbols in reverse order.

Practical implementations handle this by simply encoding the
input block \emph{right-to-left}.  Then the decoder, running
left-to-right, produces the original order automatically.
\end{importantbox}


\subsubsection{The rANS Decoding Step}
\label{sec:rans_decode}

\paragraph{What the decoder receives.}
The decoder starts from the final state $x$ (written at the
end of the bitstream) and has access to the bitstream of
pre-norm bits in \emph{reverse} order.

\paragraph{How to identify the symbol.}
Compute $r = x \bmod L$.  Look up which symbol $s$ has
$c_s \le r < c_s + f_s$.  This is the slot-range lookup
described in Section~\ref{sec:rans_encode}: the encoder
guaranteed that $x \bmod L$ always falls in $s$'s slot range,
so this lookup unambiguously identifies $s$.

\paragraph{How to recover the previous state.}
Once $s$ is known, we can invert formula \eqref{eq:rans_enc}.
Since $x \in [L, 2L)$ we have $\lfloor x/L \rfloor = 1$, so:
\[
  x_{\text{prev}} \;=\; f_s \cdot \Bigl\lfloor \frac{x}{L} \Bigr\rfloor
                  \;+\; r - c_s
                  \;=\; f_s + r - c_s
\]
Because $r \in [c_s, c_s + f_s)$, we have
$x_{\text{prev}} \in [f_s, 2f_s)$ — exactly the
pre-normalised range the encoder started from before it
encoded $s$.

\paragraph{Renormalise upward.}
$x_{\text{prev}}$ may be less than $L$, so we must restore it
to $[L, 2L)$.  We do this by reading bits from the (reversed)
bitstream and appending them to the bottom of $x_{\text{prev}}$
until $x_{\text{prev}} \ge L$.  This consumes exactly the bits
the encoder emitted for this symbol, restoring the encoder's
state to what it was \emph{before} encoding $s$.

\begin{definitionbox}
\textbf{rANS Decode Algorithm.}  Given state
$x \in [L, 2L)$, decode one symbol:

\medskip
\textbf{Step 1 — Identify symbol:}
\begin{algorithmic}[1]
\State $r \leftarrow x \bmod L$
\State Find symbol $s$ such that $c_s \le r < c_s + f_s$
       \Comment{check slot ranges; linear or binary search}
\end{algorithmic}

\medskip
\textbf{Step 2 — Recover previous state:}
\begin{algorithmic}[1]
\State $x_{\text{prev}} \leftarrow f_s + r - c_s$
       \Comment{algebraic inverse of the encode formula}
\end{algorithmic}

\medskip
\textbf{Step 3 — Renormalise upward:}
\begin{algorithmic}[1]
\While{$x_{\text{prev}} < L$}
    \State Read next bit $b$ from the (reversed) stream
    \State $x_{\text{prev}} \leftarrow (x_{\text{prev}} \ll 1) \mid b$
           \Comment{shift left and append bit at the bottom}
\EndWhile
\Comment{Now $L \le x_{\text{prev}} < 2L$; this is the new $x$}
\end{algorithmic}

\medskip
\textbf{Step 4:} Output symbol $s$.  Set $x \leftarrow x_{\text{prev}}$
and repeat from Step 1 for the next symbol.
\end{definitionbox}

% -----------------------------------------------------------
\subsection{Worked Example: rANS Encoding}
% -----------------------------------------------------------

We now trace through a complete encoding by hand so every
formula step can be verified concretely.

\begin{examplebox}
\textbf{Setup.}
\[
  \mathcal{A} = \{A, B, C\},\quad L = 8,\quad R = 3.
\]

\smallskip
\begin{center}
\begin{tabular}{@{}lcccc@{}}
\toprule
Symbol & Freq.\ $f$ & Cumul.\ $c$ & $\tilde{p}=f/L$ & Slot range $[c,\;c+f)$ \\
\midrule
$A$ & 4 & 0 & 0.50 & $[0,4)$ \\
$B$ & 2 & 4 & 0.25 & $[4,6)$ \\
$C$ & 2 & 6 & 0.25 & $[6,8)$ \\
\bottomrule
\end{tabular}
\end{center}

\medskip
\noindent\textbf{Goal:} encode the sequence \texttt{A B A},
right-to-left (so we start with the last symbol and work
backwards).  Initial state $x = L = 8$.

\medskip
The two-step procedure for each symbol:
\begin{enumerate}
  \item \textbf{Pre-normalise:} repeatedly emit the lowest
        bit of $x$ and right-shift until $x \in [f_s, 2f_s)$.
  \item \textbf{Apply formula:}
        $x \leftarrow L + c_s + (x \bmod f_s)$.
\end{enumerate}

\medskip\hrule\medskip

\noindent\textbf{Step 1 --- Encode the 3rd symbol: \texttt{A}}

\noindent\emph{Parameters:} $f_A = 4$, $c_A = 0$.
Current state: $x = 8$.

\noindent\emph{Pre-normalise into $[f_A, 2f_A) = [4, 8)$:}
\begin{algorithmic}[1]
\While{$x \ge 2f_A = 8$}
    \State $x = 8 \ge 8$: \enspace emit $(8 \mathrel{\&} 1) = \mathbf{0}$;\enspace
           $x \leftarrow 8 \gg 1 = 4$
\EndWhile
\Comment{$4 \in [4,8)$: stop}
\end{algorithmic}

\noindent\emph{Apply formula} with pre-normed $x = 4$:
\[
  x \leftarrow 8 + 0 + (4 \bmod 4) = 8 + 0 + 0 = 8
\]

\begin{center}
\begin{tabular}{@{}ll@{}}
  Bit emitted:      & \texttt{0} \\
  Bitstream so far: & \texttt{0} \\
  New state $x$:    & $8$ \\
\end{tabular}
\end{center}

\noindent\textit{Why 1 bit?}  $A$ has probability 0.5, so it
carries $-\log_2 0.5 = 1$ bit of information.  Pre-normalising
from $x=8$ to $[4,8)$ required one right-shift, emitting 1
bit.

\medskip\hrule\medskip

\noindent\textbf{Step 2 --- Encode the 2nd symbol: \texttt{B}}

\noindent\emph{Parameters:} $f_B = 2$, $c_B = 4$.
Current state: $x = 8$.

\noindent\emph{Pre-normalise into $[f_B, 2f_B) = [2, 4)$:}
\begin{algorithmic}[1]
\While{$x \ge 2f_B = 4$}
    \State $x = 8 \ge 4$: \enspace emit $(8 \mathrel{\&} 1) = \mathbf{0}$;\enspace
           $x \leftarrow 8 \gg 1 = 4$
    \State $x = 4 \ge 4$: \enspace emit $(4 \mathrel{\&} 1) = \mathbf{0}$;\enspace
           $x \leftarrow 4 \gg 1 = 2$
\EndWhile
\Comment{$2 \in [2,4)$: stop}
\end{algorithmic}

\noindent\emph{Apply formula} with pre-normed $x = 2$:
\[
  x \leftarrow 8 + 4 + (2 \bmod 2) = 8 + 4 + 0 = 12
\]

\begin{center}
\begin{tabular}{@{}ll@{}}
  Bits emitted:     & \texttt{0\;0} \\
  Bitstream so far: & \texttt{0\;0\;0} \\
  New state $x$:    & $12$ \\
\end{tabular}
\end{center}

\noindent\textit{Why 2 bits?}  $B$ has probability 0.25, so
it carries $-\log_2 0.25 = 2$ bits.  We needed two
right-shifts to bring $x = 8$ into $[2, 4)$.

\medskip\hrule\medskip

\noindent\textbf{Step 3 --- Encode the 1st symbol: \texttt{A}}

\noindent\emph{Parameters:} $f_A = 4$, $c_A = 0$.
Current state: $x = 12$.

\noindent\emph{Pre-normalise into $[f_A, 2f_A) = [4, 8)$:}
\begin{algorithmic}[1]
\While{$x \ge 2f_A = 8$}
    \State $x = 12 \ge 8$: \enspace emit $(12 \mathrel{\&} 1) = \mathbf{0}$;\enspace
           $x \leftarrow 12 \gg 1 = 6$
\EndWhile
\Comment{$6 \in [4,8)$: stop}
\end{algorithmic}

\noindent\emph{Apply formula} with pre-normed $x = 6$:
\[
  x \leftarrow 8 + 0 + (6 \bmod 4) = 8 + 0 + 2 = 10
\]

\begin{center}
\begin{tabular}{@{}ll@{}}
  Bit emitted:      & \texttt{0} \\
  Bitstream so far: & \texttt{0\;0\;0\;0} \\
  New state $x$:    & $10$ \\
\end{tabular}
\end{center}

\medskip\hrule\medskip

\noindent\textbf{Finalise.}  Append the final state
$x = 10 = (1010)_2$ to the bitstream:

\[
\underbrace{\mathtt{0\;0\;0\;0}}_{\text{pre-norm bits from steps 1--3}}
\;\Big|\;
\underbrace{\mathtt{1\;0\;1\;0}}_{\text{final state } x = 10}
\]

\medskip
\noindent\textbf{Complete encoding summary:}

\smallskip
\begin{center}
\begin{tabular}{@{}cclccc@{}}
\toprule
Step & Symbol & $x$ in & Bits emitted & $x$ after pre-norm & $x$ out \\
\midrule
1 & $A$ & 8  & \texttt{0}    & 4 & 8  \\
2 & $B$ & 8  & \texttt{0\;0} & 2 & 12 \\
3 & $A$ & 12 & \texttt{0}    & 6 & 10 \\
\bottomrule
\end{tabular}
\end{center}

\medskip
\noindent\textit{Entropy check.}  Shannon entropy of
\texttt{A B A}: $-(2 \times \log_2 0.5 + \log_2 0.25)
= 1 + 1 + 2 = 4$ bits.  We emitted exactly 4 pre-norm bits.
The additional 4 bits storing the final state $x=10$ are a
fixed per-block overhead that becomes negligible for long
sequences.
\end{examplebox}


% -----------------------------------------------------------
\subsection{Worked Example: rANS Decoding}
% -----------------------------------------------------------

We now decode the bitstream produced in the previous encoding
example, step by step, verifying that we recover \texttt{A B A}.

\begin{examplebox}
\textbf{What the decoder knows:}
\begin{itemize}
  \item The final state written by the encoder: $x = 10$ (binary $1010_2$).
  \item The pre-normalisation bits emitted during encoding:
        \texttt{0, 0, 0, 0} (in the order they were emitted).
  \item The symbol table (same as encoding):
        \[
        A \to [0,4),\qquad B \to [4,6),\qquad C \to [6,8)
        \]
  \item $L = 8$, $R = 3$.
\end{itemize}

\medskip
\textbf{Critical note:} The decoder reads pre-normalisation bits
in \textbf{reverse order} (because the encoder emitted bits for the
last symbol first).  In this example all bits are zero, so reversal
doesn't change them, but the principle matters.

\begin{center}
\begin{tabular}{@{}ll@{}}
  Initial state $x$:  & $10$ \\
  Bit queue (reversed): & \texttt{0, 0, 0, 0} (read left-to-right) \\
\end{tabular}
\end{center}

\medskip\hrule\medskip

\noindent\textbf{Decode Step 1 — Recover the last encoded symbol (position 3)}:

\medskip
\noindent\emph{Step 1a: Identify the symbol.}
\[
r = x \bmod L = 10 \bmod 8 = 2
\]
Check which symbol owns slot $2$:
\begin{itemize}
  \item $A$ owns $[0,4)$ → $2$ is in $A$'s range.
  \item $B$ owns $[4,6)$ → no.
  \item $C$ owns $[6,8)$ → no.
\end{itemize}
Thus $s = A$.

\medskip
\noindent\emph{Step 1b: Recover the previous state before encoding this $A$.}
\[
x_{\text{prev}} = f_A + r - c_A = 4 + 2 - 0 = 6
\]

\noindent\emph{Step 1c: Renormalise upward.}
The encoder had pre-normalised $x$ into $[f_A, 2f_A) = [4,8)$
before encoding.  Our $x_{\text{prev}} = 6$ is already in that range,
but it may be below $L = 8$.  We need to restore it to $[L, 2L) = [8,16)$
by reading bits from the stream and appending them to the bottom:

\begin{algorithmic}[1]
\While{$x_{\text{prev}} < 8$}
    \State $6 < 8$: read next bit $b = \mathbf{0}$
    \State $x_{\text{prev}} \leftarrow (6 \ll 1) \mid 0 = 12$
\EndWhile
\Comment{$12 \ge 8$: stop}
\end{algorithmic}

\begin{center}
\begin{tabular}{@{}ll@{}}
  Decoded symbol: & $A$ \\
  New state $x$:  & $12$ \\
  Bits remaining: & \texttt{0, 0, 0} \\
\end{tabular}
\end{center}

\noindent\textit{Why this works:} The encoder had emitted bits
until $x$ fell into $[4,8)$.  By reading those same bits back
in reverse (here just one zero) and shifting them in from below,
we reconstruct the encoder's state \emph{before} it started
pre-normalising.

\medskip\hrule\medskip

\noindent\textbf{Decode Step 2 — Recover the second-last encoded symbol (position 2)}:

\medskip
\noindent\emph{Step 2a: Identify the symbol.}
\[
r = x \bmod L = 12 \bmod 8 = 4
\]
Slot $4$ belongs to $B$ (since $B$ owns $[4,6)$).  Thus $s = B$.

\medskip
\noindent\emph{Step 2b: Recover the previous state.}
\[
x_{\text{prev}} = f_B + r - c_B = 2 + 4 - 4 = 2
\]

\noindent\emph{Step 2c: Renormalise upward.}
We need $x_{\text{prev}} \ge L = 8$:

\begin{algorithmic}[1]
\While{$x_{\text{prev}} < 8$}
    \State $2 < 8$: read next bit $b = \mathbf{0}$
    \State $x_{\text{prev}} \leftarrow (2 \ll 1) \mid 0 = 4$
    \State $4 < 8$: read next bit $b = \mathbf{0}$
    \State $x_{\text{prev}} \leftarrow (4 \ll 1) \mid 0 = 8$
\EndWhile
\Comment{$8 \ge 8$: stop}
\end{algorithmic}

\begin{center}
\begin{tabular}{@{}ll@{}}
  Decoded symbol: & $B$ \\
  New state $x$:  & $8$ \\
  Bits remaining: & \texttt{0} \\
\end{tabular}
\end{center}

\noindent\textit{Why 2 bits?}  $B$ has probability $0.25$,
so it should add $-\log_2(0.25)=2$ bits of information.
The encoder emitted exactly 2 bits for this symbol, and the
decoder now consumes them.

\medskip\hrule\medskip

\noindent\textbf{Decode Step 3 — Recover the first encoded symbol (position 1)}:

\medskip
\noindent\emph{Step 3a: Identify the symbol.}
\[
r = x \bmod L = 8 \bmod 8 = 0
\]
Slot $0$ belongs to $A$.  Thus $s = A$.

\medskip
\noindent\emph{Step 3b: Recover the previous state.}
\[
x_{\text{prev}} = f_A + r - c_A = 4 + 0 - 0 = 4
\]

\noindent\emph{Step 3c: Renormalise upward.}

\begin{algorithmic}[1]
\While{$x_{\text{prev}} < 8$}
    \State $4 < 8$: read next bit $b = \mathbf{0}$
    \State $x_{\text{prev}} \leftarrow (4 \ll 1) \mid 0 = 8$
\EndWhile
\Comment{$8 \ge 8$: stop}
\end{algorithmic}

\begin{center}
\begin{tabular}{@{}ll@{}}
  Decoded symbol: & $A$ \\
  New state $x$:  & $8$ (the encoder's initial state $L$) \\
  Bits remaining: & (empty) \\
\end{tabular}
\end{center}

\medskip\hrule\medskip

\noindent\textbf{Decoding complete.}  The decoded symbols in order
(by the order we output them) are:
\[
\text{Step 1: } A,\quad \text{Step 2: } B,\quad \text{Step 3: } A
\]
Thus we recover \texttt{A B A} — identical to the input. \checkmark

\medskip
\noindent\textbf{Summary table for all decoding steps:}

\begin{center}
\begin{tabular}{@{}cclclcl@{}}
\toprule
Step & $x$ (in) & $r = x \bmod 8$ & Symbol & $x_{\text{prev}}$ (raw) & Bits read & $x$ (out) \\
\midrule
1 & 10 & 2 & $A$ & 6 & \texttt{0}    & 12 \\
2 & 12 & 4 & $B$ & 2 & \texttt{0\;0} & 8  \\
3 & 8  & 0 & $A$ & 4 & \texttt{0}    & 8  \\
\bottomrule
\end{tabular}
\end{center}

\medskip
\noindent\textbf{Verification of information content:}
\begin{itemize}
  \item The decoder consumed exactly 4 bits from the stream.
  \item Shannon entropy of \texttt{A B A} is $1 + 1 + 2 = 4$ bits.
  \item No bits are wasted — the decoder recovers the exact
        original symbols with perfect efficiency.
\end{itemize}
\end{examplebox}


% -----------------------------------------------------------
\subsection{Why the Encoding Formula Works: A Visual Summary}
% -----------------------------------------------------------

The encoding formula $x' = L + c_s + (x \bmod f_s)$ does
two things simultaneously:

\begin{center}
\renewcommand{\arraystretch}{1.4}
\begin{tabular}{@{}p{4cm}p{4.5cm}p{5cm}@{}}
\toprule
\textbf{Part of formula} & \textbf{What it does} & \textbf{Why} \\
\midrule
$L$ \newline (from $\lfloor x/f_s\rfloor \cdot L$
after pre-norm) &
Adds $\log_2(L/f_s)$ bits to the state &
This equals $-\log_2(f_s/L) = -\log_2 \tilde{p}_s$, the
ideal bit-cost of symbol $s$ \\
\midrule
$c_s + (x \bmod f_s)$ &
Plants a ``fingerprint'' in the low $R$ bits of $x'$ &
Ensures $x' \bmod L \in [c_s, c_s+f_s)$, so the decoder
can identify $s$ by a single range check \\
\bottomrule
\end{tabular}
\end{center}

The elegance of ANS is that both jobs are done by a single
addition and a single modulo — no separate tagging step is
needed.

% -----------------------------------------------------------
\subsection{tANS: Tabled ANS}
\label{sec:tans}
% -----------------------------------------------------------

\subsubsection{Motivation: Eliminating Division}

The rANS pre-normalisation loop divides $x$ by 2 repeatedly
until $x \in [f_s, 2f_s)$.  On each loop iteration we
perform the check $x \ge 2f_s$ and a right-shift.  The
number of iterations differs per symbol, making the loop
hard to pipeline on modern CPUs.

\textbf{tANS eliminates all of this} by precomputing, for
every possible (symbol, input state) pair, the exact output
state and the number of bits to emit.  At runtime the entire
encode or decode step is a \emph{single array lookup}.

\subsubsection{Step 1: Spreading Symbols into Slots}
\label{sec:ans_spread}

Recall that we have $L$ slots numbered $0$ to $L-1$, and
symbol $s$ must occupy exactly $f_s$ of them.  For the
lookup table to work well, we want each symbol's copies to
be spread as \emph{evenly} as possible across the slots
(rather than grouped together), so that consecutive states
visit a balanced mix of symbols.

The even-spacing formula (used by Zstandard) places the
$i$-th copy of symbol $s$ at slot:
\[
  \text{slot}(s, i) =
  \Bigl( c_s + \Bigl\lfloor \frac{i \cdot L}{f_s} \Bigr\rfloor
  \Bigr) \bmod L,
  \qquad i = 0, 1, \ldots, f_s - 1.
\]
The result is an array \texttt{sym[0..L-1]} where each entry
names the symbol that owns that slot.

\subsubsection{Step 2: Building the Decode Table}

Once \texttt{sym[]} is fixed, for each slot $q$ we
precompute the previous encoder state:
\[
  x_{\text{prev}}(q)
  = f_s + (q \bmod f_s) - c_s,
  \qquad s = \texttt{sym}[q].
\]
This is exactly the rANS decode formula
$x_{\text{prev}} = f_s + r - c_s$ with $r = (L+q) \bmod L = q$,
evaluated for the ``canonical'' state $x = L + q$
(i.e.\ $\lfloor x/L\rfloor = 1$).

If $x_{\text{prev}}(q) < L$, bits must be read to bring it
into $[L, 2L)$.  A real implementation stores the
\emph{post-renorm} value directly so that no computation
happens at decode time.

\subsubsection{Step 3: Building the Encode Table}

For each symbol $s$ and each possible value $v = x \bmod f_s
\in \{0, 1, \ldots, f_s-1\}$, the encode table stores:
\begin{itemize}
  \item The number of bits to emit (how many right-shifts
        are needed to bring $x$ into $[f_s, 2f_s)$).
  \item The resulting new state after applying the formula.
\end{itemize}
At runtime: look up \texttt{enc[s][$x \bmod f_s$]}, emit
the stored bits, and set $x$ to the stored new state.
No loop, no division.

\subsubsection{Full tANS Example}

\begin{examplebox}
\textbf{Setup.}
$\mathcal{A} = \{A, B, C\}$, $L = 8$;
$f_A = 4,\; c_A = 0$;\;
$f_B = 2,\; c_B = 4$;\;
$f_C = 2,\; c_C = 6$.

\medskip
\textbf{1.\ Spread symbols.}
Apply the even-spacing formula for each symbol:

\begin{center}
\begin{tabular}{@{}ccccccccc@{}}
\toprule
Slot $q$        & 0 & 1 & 2 & 3 & 4 & 5 & 6 & 7 \\
\midrule
\texttt{sym[q]} & A & C & B & A & C & B & A & A \\
\bottomrule
\end{tabular}
\end{center}

$A$ occupies slots $\{0,3,6,7\}$;\quad
$B$ occupies slots $\{2,5\}$;\quad
$C$ occupies slots $\{1,4\}$.

\medskip
\textbf{2.\ Build the decode table.}

For each slot $q$, look up $s = \texttt{sym}[q]$ and compute
$x_{\text{prev}}(q) = f_s + (q \bmod f_s) - c_s$:

\begin{center}
\begin{tabular}{@{}ccccll@{}}
\toprule
$q$ & $s$ & $f_s$ & $c_s$ & Computation & $x_{\text{prev}}$ \\
\midrule
0 & $A$ & 4 & 0 & $4+(0\bmod 4)-0$   & $4$ \quad (renorm: $4<8$, read bits) \\
1 & $C$ & 2 & 6 & $2+(1\bmod 2)-6$   & $-3$ \quad (renorm needed) \\
2 & $B$ & 2 & 4 & $2+(2\bmod 2)-4$   & $-2$ \quad (renorm needed) \\
3 & $A$ & 4 & 0 & $4+(3\bmod 4)-0$   & $7$ \quad ($7<8$, read 1 bit) \\
4 & $C$ & 2 & 6 & $2+(4\bmod 2)-6$   & $-4$ \quad (renorm needed) \\
5 & $B$ & 2 & 4 & $2+(5\bmod 2)-4$   & $-1$ \quad (renorm needed) \\
6 & $A$ & 4 & 0 & $4+(6\bmod 4)-0$   & $6$ \quad (renorm: read bits) \\
7 & $A$ & 4 & 0 & $4+(7\bmod 4)-0$   & $7$ \quad (renorm: read 1 bit) \\
\bottomrule
\end{tabular}
\end{center}

All raw $x_{\text{prev}}$ values below $L=8$ must be
renormalised upward by reading bits.  The table stores the
final post-renorm values so decode is a single lookup.

\medskip
\textbf{3.\ Decode one step.}

Suppose the final encoder state is $x = 8$.
Compute $q = x \bmod L = 0$.

\begin{itemize}
  \item \texttt{sym[0]} $= A$, $x_{\text{prev}}(0) = 4$.
  \item $4 < 8$: read one bit $b$.
        \begin{itemize}
          \item If $b=0$: $x_{\text{prev}} = (4 \ll 1)\mid 0 = 8$.
          \item If $b=1$: $x_{\text{prev}} = (4 \ll 1)\mid 1 = 9$.
        \end{itemize}
        Either way $x_{\text{prev}} \ge 8$: done.
  \item Output symbol $A$.  Decoded in one lookup + one
        bit-read.  No arithmetic.
\end{itemize}
\end{examplebox}

% -----------------------------------------------------------
\subsection{ANS vs.\ Arithmetic Coding vs.\ Huffman}
% -----------------------------------------------------------

\begin{center}
\rowcolors{2}{blue!8}{white}
\begin{tabular}{@{}p{3.5cm}p{3.5cm}p{3.5cm}p{3.5cm}@{}}
\toprule
\textbf{Property} & \textbf{Huffman}
  & \textbf{Arithmetic} & \textbf{ANS (rANS/tANS)} \\
\midrule
Compression ratio & Near-entropy (integer codes)
  & Near-entropy & Near-entropy \\
Speed (encode)    & Very fast  & Moderate & Fast \\
Speed (decode)    & Very fast  & Moderate & Very fast (tANS) \\
Division at runtime? & No & Yes & No (tANS) / Yes (rANS) \\
Parallelisable?   & Limited    & Hard     & Easy (interleaved) \\
Decode order      & Forward    & Forward  & Reverse (auto-handled) \\
Patent issues     & None       & Some (old) & None \\
\bottomrule
\end{tabular}
\end{center}

% -----------------------------------------------------------
\subsection{Why ANS Achieves Near-Entropy Compression}
% -----------------------------------------------------------

\begin{theorem}
Let $p_s$ be the true probability of symbol $s$ and
$\tilde{p}_s = f_s / L$ the quantised probability.  The
redundancy (extra bits per symbol above Shannon entropy) is
bounded by:
\[
  \Delta H \;\leq\;
  \underbrace{\sum_s p_s \log_2 \frac{p_s}{\tilde{p}_s}}_{\text{quantisation error}}
  \;+\;
  \underbrace{\frac{1}{L \ln 2}}_{\text{rounding error}}
\]
\end{theorem}

\begin{itemize}[leftmargin=2em]
  \item The \textbf{quantisation error} is the KL divergence
        between the true and quantised distributions.  It
        vanishes as $L \to \infty$ because $\tilde{p}_s \to p_s$.

  \item The \textbf{rounding error} $1/(L\ln 2)$ is tiny
        even for moderate $L$: with $L = 4096$ it is less
        than $0.0004$ bits per symbol.

  \item In practice, $L = 2^{11}$ or $2^{12}$ gives
        compression within $0.001$ bits per symbol of the
        Shannon limit — far tighter than Huffman coding.
\end{itemize}

% -----------------------------------------------------------
\subsection{Interleaved ANS for Parallelism}
% -----------------------------------------------------------

A powerful extension of ANS is to maintain \emph{multiple}
independent state variables $x_0, x_1, \ldots$ that each
encode a separate sub-stream.  Because the states are
independent, their encode/decode steps can be executed in
parallel using SIMD (Single Instruction Multiple Data) CPU
instructions.

\begin{examplebox}
\textbf{2-way interleaved rANS.}
Maintain two states $x_0$ and $x_1$.  Encode even-indexed
symbols (positions 0, 2, 4, \ldots) into $x_0$ and
odd-indexed symbols (positions 1, 3, 5, \ldots) into $x_1$.
The two bitstreams are concatenated for storage.

Zstandard uses \textbf{4-way interleaving}, allowing it to
process four symbols simultaneously and achieving decompression
throughput of several gigabytes per second on modern hardware.
\end{examplebox}

% -----------------------------------------------------------
\subsection{ANS in Practice: Zstandard}
% -----------------------------------------------------------

Zstandard (zstd), developed at Meta and standardised as
RFC~8878, uses ANS as its entropy back-end:

\begin{itemize}[leftmargin=2em]
  \item \textbf{Literals} (raw byte values) are compressed
        with Huffman or FSE (Finite State Entropy — Meta's
        name for tANS).
  \item \textbf{Offsets, match lengths, literal lengths}
        (the LZ77 back-reference parameters) are compressed
        with FSE tables.
  \item The table size $L$ ranges from $2^6$ to $2^9$,
        chosen adaptively for each block.
  \item At compression level 3 (the default), zstd
        decompresses roughly $5\times$ faster than gzip
        while achieving similar or better compression ratios.
\end{itemize}

% -----------------------------------------------------------
\subsection{Summary of Key Concepts}
% -----------------------------------------------------------

\begin{center}
\rowcolors{2}{green!8}{white}
\begin{tabular}{@{}p{4.5cm}p{9.5cm}@{}}
\toprule
\textbf{Concept} & \textbf{Plain-English Summary} \\
\midrule
State $x$ &
  A single integer accumulating all encoded information.
  Its binary size tracks the entropy of the input. \\
Normalised range $[L, 2L)$ &
  Keeps $x$ from growing unboundedly.  Bits are emitted
  (encoder) or consumed (decoder) to maintain this range. \\
Frequency $f_s$ &
  Integer approximation to $p_s \cdot L$.
  Determines how many slots symbol $s$ owns. \\
Cumulative frequency $c_s$ &
  Starting slot for symbol $s$: $c_s = \sum_{j<s} f_j$.
  Used to convert between slot numbers and symbols. \\
Slot range $[c_s, c_s+f_s)$ &
  The block of slots owned by $s$.
  Decoder identifies $s$ by checking which range $x \bmod L$ falls in. \\
Pre-normalisation &
  Emit bits from $x$ until $x \in [f_s, 2f_s)$, \emph{before}
  applying the encode formula, so the output stays in $[L,2L)$. \\
Encode formula &
  $x \leftarrow L + c_s + (x \bmod f_s)$ (after pre-norm).
  Adds $-\log_2 \tilde{p}_s$ bits and tags the result for the decoder. \\
Decode formula &
  $x_{\text{prev}} = f_s + (x \bmod L) - c_s$.
  Exact algebraic inverse of the encode formula. \\
tANS &
  Precomputes all encode/decode arithmetic into two arrays.
  Reduces each step to one table lookup plus bit shifts. \\
Symbol spread &
  Places $f_s$ copies of symbol $s$ evenly across the $L$
  slots so that consecutive states visit a balanced mix. \\
Encode backwards &
  Encode the input right-to-left so that decoding
  (which runs in reverse) produces the original left-to-right order. \\
Redundancy &
  Approaches 0 as table size $L \to \infty$;
  already negligible ($<0.001$ bits/symbol) for $L = 2^{12}$. \\
\bottomrule
\end{tabular}
\end{center}

% -----------------------------------------------------------
\subsection{Exercises}
% -----------------------------------------------------------

\begin{exercisebox}
\label{ex:ans1}
Consider the alphabet $\{X, Y\}$ with frequencies $f_X = 3$,
$f_Y = 1$, and table size $L = 4$.
\begin{enumerate}[label=(\alph*)]
  \item Compute the cumulative frequencies $c_X$ and $c_Y$.
        Write out the slot ranges for $X$ and $Y$.
  \item Using the rANS encode procedure (pre-normalise into
        $[f_s, 2f_s)$, then apply $x \leftarrow L + c_s +
        (x \bmod f_s)$), encode the sequence \texttt{X X Y}
        starting from state $x = 4$.  Remember to encode
        right-to-left (Y first, then X, then X).  Show every
        intermediate state and every bit emitted.
  \item Decode the result step by step and verify you recover
        \texttt{X X Y}.
\end{enumerate}
\end{exercisebox}

\begin{exercisebox}
\textbf{Entropy comparison.}  For the alphabet in
Exercise~\ref{ex:ans1}:
\begin{enumerate}[label=(\alph*)]
  \item Compute the Shannon entropy $H = -\sum_s p_s \log_2 p_s$.
  \item Count the total bits produced by your rANS encoding.
  \item Compute the redundancy as a percentage above $H$.
  \item How would increasing $L$ to 16 affect the redundancy,
        and why?
\end{enumerate}
\end{exercisebox}

\begin{exercisebox}
\textbf{tANS table construction.}
Given $\mathcal{A} = \{A, B\}$, $L = 8$, $f_A = 5$,
$f_B = 3$:
\begin{enumerate}[label=(\alph*)]
  \item Use the even-spacing formula from
        Section~\ref{sec:ans_spread} to spread symbols and
        write the \texttt{sym[0..7]} array.
  \item Build the decode table: for each slot
        $q \in \{0,\ldots,7\}$, compute the symbol and the
        raw $x_{\text{prev}}$.  Note which entries require
        upward renormalisation.
  \item Starting from state $x = 10$, perform one decode step:
        find the symbol, compute $x_{\text{prev}}$, and
        renormalise if necessary.
\end{enumerate}
\end{exercisebox}

\begin{exercisebox}
\textbf{Conceptual questions.}
\begin{enumerate}[label=(\alph*)]
  \item Explain in your own words why the encoder must
        pre-normalise $x$ into $[f_s, 2f_s)$ \emph{before}
        applying the encoding formula.  What would go wrong
        if we applied the formula first and normalised
        afterwards?
  \item Why does encoding a rare symbol (small $f_s$) emit
        more bits during pre-normalisation than encoding a
        common symbol?  Relate your answer to the entropy
        formula $-\log_2 p_s$.
  \item Why does ANS decode in reverse order?  What do
        practical implementations do to obtain the correct
        output order?
  \item What is the advantage of interleaved ANS over a
        single-state implementation on modern CPUs?
  \item Name two software products that use ANS and state
        which variant (rANS or tANS) each employs.
\end{enumerate}
\end{exercisebox}
